%======================================================================
\section{Discussion and limitations}
\label{sec:discussion}
%======================================================================

\para{What transfer requires.} Four conditions emerged, and each was
necessary in at least one system. First, the physics must be extensive
at the scales involved, which can be checked from ground truth at two
sizes before any model is trained: the 2D KS density is constant within
the spread of the truth runs ($82.0$ and $84.0$ per base box at $44^2$ and
$66^2$), while in Kolmogorov flow the truth's density per box rises with
the channel width at every drag (Table~\ref{tab:truth}). A local model
is extensive by construction (Remark~\ref{rem:extensive_model}), so
wherever the truth is not, a transfer error is expected and is not a
failure of training. Second, the training data must contain the local
statistics of a large domain: trained on a small periodic box the same
architecture misses $\DKY$ by $15$--$38\%$ in three of four larger
configurations, and trained on windows of one large simulation it
transfers (Sec.~\ref{sec:box_vs_window}); the window training is a
limited-area problem, and Proposition~\ref{prop:halo} shows that the
forced halo trains exactly that problem. Third, the learned map must
remain chaotic, and chaotic at the right rate, when it runs freely
(Secs.~\ref{sec:fail_regimes} and \ref{sec:fail_mse}); free-rollout
statistics detect spurious regimes without any truth, and the truth-free
selection rule of Sec.~\ref{sec:canonical_results} chose a model within
$0.9$--$5.9\%$ of the truth at every width. Fourth, the model's reach must
cover the interaction range of the dynamics.

\para{The interaction range is the design length.} The light cone is far
too short to constrain the reach of a model at these time steps; the
length that matters is the distance beyond which neighbouring values stop
improving a one-step prediction, three to four sites for the 1D latent
dynamics against a trained reach of six (Sec.~\ref{sec:receptive}). In
Kolmogorov flow most of the unstable directions lie in the band around
the forcing wavenumber, $31.5$ of $36$ at $|\bm k|\in[2,8)$ for $\alpha=0.05$ at
$3\times$, whose wavelengths span one to four sites, and a model of finite
reach reproduces this scale content (Appendix~\ref{app:chaos_scale}). Models whose encoder
reaches two sites and whose propagator reaches one fail even when trained
on windows of the very domain they are evaluated on, with errors of
$66$--$100\%$ in $\DKY$ (Appendix~\ref{app:tables}). The rule we take from
this is to measure the interaction range on data, for example by fitting
a shared stencil of increasing width, and to give the encoder and
propagator a reach that covers it with a margin. The reach then sets the
smallest domain the model can be used on, which is what fails at $L=25$.

\para{The parameters describe the law, not the state.} Because every
operator is shared by all sites, the trained parameters encode one local
rule, and their number is fixed by that rule rather than by the domain.
A 2D model of about $135{,}000$ parameters reproduces a Kaplan--Yorke
dimension of $188$ at $66^2$ to within $0.6\%$ (Table~\ref{tab:transfer2d}), and
a 12-channel version runs a $440^2$ domain, whose attractor dimension we
estimate at roughly $8{,}000$ from the measured density of $0.043$ per unit area
\citep{jones2026da}. The same weight sharing makes training cheap in
data: every window, and every site within it, is a training example of
the same rule. Weight sharing and size independence are not new by
themselves \citep{barsinai2019,kochkov2021}; what the dimension comparison
adds is evidence that the shared rule is the right one, in the sense that
its extension to a larger domain carries the right density of chaos. A
system made of many copies of one local process is described most
economically by learning that process once.

\para{Transfer error as a probe of the physics.} Because a local model
is extensive by construction, the difference between its dimension
density and the truth's on a larger domain measures how far the truth is
from extensive at that size, once the model is accurate at the native
size. The weakest drag, where the native-size models scatter most and the
transferred models fall below a rising truth, is where the inverse
cascade is least arrested, and we expect non-local effects to be largest
there. Using
transfer error to locate the scales at which a system stops being
extensive is a question in its own right, which we leave to future work.

\para{Limitations.} We state them plainly.
\begin{itemize}
\item \emph{Ground truth.} Five Kolmogorov configurations rest on one 40-unit
  run, uncertain by 10--15\%, the 2D KS truths by a few percent
  (Sec.~\ref{sec:truth_uncertainty}), and there is no truth at $88^2$.
\item \emph{Kolmogorov geometry.} The channel is widened in $x$ only,
  because the forcing fixes the height; transfer in $y$ is untested.
\item \emph{Models.} Most 2D models have one training seed (the 1D model
  has three, the $\alpha=0.05$ canonical model two); the 1D and 2D
  canonical architectures differ (full reach in 1D, per-patch in 2D); the
  weakest drag remains $8$--$25\%$ off.
\item \emph{What the spectra compare.} The model's Lyapunov spectrum is
  that of the latent map. Comparing it with the PDE's assumes that the
  decoder is a smooth, injective map on the attractor, which the agreement
  of the spectra supports but does not prove.
\item \emph{Reconstruction.} Local neural-field latents have been
  reported to fail on KS at the decoder \citep{aroma2024}. Ours reconstruct
  the field to $0.01\%$ (1D), $2.6$--$2.9\%$ (2D KS) and $1.2\%$ (Kolmogorov) of
  its variance, with seams at patch edges in 2D.
\end{itemize}

\para{Outlook.} The ingredients carry over to any system whose chaos is
extensive and whose dynamics are local at the scale of a few patches:
convection in large aspect-ratio cells, limited-area atmospheric models
nested in a global one, where the forced halo is already the operating
mode, and high-resolution models whose domain is chosen by the
application rather than by the training data. Three questions follow
from this study. Can window training be made to work at the weakest drag,
where the inverse cascade is least arrested? Does transfer hold in a
direction the training data never varied, such as the height of the
Kolmogorov channel? And can the rule be learned from incomplete, noisy
data rather than from complete fields? A companion paper uses the local
latent field for ensemble data assimilation, with a model trained only on
noisy fields \citep{jones2026da}.
